%======================================================================
\section{Discussion}\label{sec:discussion}
%======================================================================

\subsection{What the AT1-specific premium is}

The AT1-specific hazard, $0.4$--$0.5\%$ a year in the base case, about $50$bp of spread, is
the part of the AT1's compensation that BNP's Tier~2 does not earn and that the model's MDA
rule does not explain. Proposition~\ref{prop:clocks} gives it a
precise risk-neutral meaning: it is the intensity of write-downs that hit the AT1 without a
resolution that also hits Tier~2, plus whatever premium investors attach to the instrument
beyond its expected loss. Four components are plausible, and the model cannot separate them.
\begin{enumerate}
  \item \emph{Write-down without resolution.} The Credit Suisse case showed that AT1 holders
  can lose everything in a transaction that leaves more junior or more senior claims
  partially intact. The EU authorities' March 2023 statement makes this less likely in the
  EU, but it does not make it impossible, and investors price the tail.
  \item \emph{Coupon discretion beyond the MDA.} The issuer may cancel coupons for reasons
  other than the MDA, and supervisors can restrict distributions outside the MDA rule, as
  they did with dividends in 2020. The model's MDA rule does not capture this.
  \item \emph{Equity-like risk premium.} An AT1 bond is perpetual, its payoff in bad states is
  close to that of equity, and its investor base includes funds with mandates that force
  sales in stress. The premium for bearing this risk need not be the same as Tier~2's.
  \item \emph{Liquidity.} AT1 bonds trade with wider bid--offer spreads than Tier~2 and are
  excluded from some benchmark indices.
\end{enumerate}
Because the MDA component is an upper bound, the $50$bp is a lower bound for these four
together, given the speed of capital mean reversion. It is not a lower bound across values
of that speed. With slower capital rebuilding, the model assigns more of the AT1's discount
to expected MDA coupon cuts and to extension and less to these four components, and at the
historical estimate it assigns them nothing (Table~\ref{tab:decompkappa}). What the data
identify is the total that the AT1 pays over the Tier~2-priced risk, $89$--$98$bp, which
does not move with the capital dynamics; its split rests on the capital dynamics and on
the MDA rule, the two weakest parts of the model.

\subsection{Implications}

\paragraph{For investors.} Two-thirds of an AT1's spread is compensation that the same
issuer's Tier~2 also pays, at a better recovery and without extension or coupon risk. The
relevant relative-value question is therefore whether the remaining $89$--$98$bp, the
AT1-specific premium together with the contract features, is adequate compensation for
those features, not whether $275$bp is a high spread in absolute terms. That total is
better determined than its parts. The model provides a consistent way to
make that comparison across issuers: calibrate $\lres$ to each issuer's Tier~2 and compare the
residuals.

\paragraph{For risk management.} The price is flat in capital at today's level and steep
within a point of the MDA threshold. A risk system that measures the distance to the trigger
will understate the capital risk of these bonds; the relevant distance is to the threshold,
which itself moves with scheduled buffer changes. The capital delta at the threshold is
three to eight times the delta today, depending on the speed of capital mean reversion
(Table~\ref{tab:kappa}).

\paragraph{For the extension decision.} The low-margin bond is the most exposed to extension
under an economic call rule, and the market behaved accordingly in September 2026. Whether
the issuer would in fact extend is a question about its reputational incentives; the
robustness results show that this question affects the extension component of the spread,
but not the hazard or its split.

\subsection{Limitations}

\begin{enumerate}
  \item \textbf{Measure.} CET1 dynamics are estimated under the physical measure and assumed
  to hold under the pricing measure. All residual premia are in $\lambda$, which is a
  risk-neutral intensity. The Tier~2 anchor splits it, but $\lres$ is risk-neutral too and
  contains Tier~2's own premia.
  \item \textbf{Recovery.} The split depends on the Tier~2 recovery (Table~\ref{tab:recovery}).
  The market prices are consistent with the conventional $20\%$, but recovery is not observed.
  \item \textbf{Capital mean reversion.} $\kappa$ is weakly identified by history, and the
  capital stress results scale by a factor of five across its bootstrap interval. The
  Tier~2 restriction $\lspec\ge0$ excludes the lower half of that interval and reduces the
  factor to about two, but it is itself conditional on the stylised MDA rule. The
  AT1-specific premium depends on $\kappa$; the Tier~2-priced hazard does not.
  \item \textbf{MDA mechanics.} Applying the quartile factor to the coupon overstates the MDA
  cost. The leverage-ratio and MREL MDAs, the available distributable items and supervisory
  discretion over distributions are not modelled.
  \item \textbf{Call behaviour.} The base rule is economic; a reputational tolerance changes
  the call probabilities but not the hazard. The regulatory approval test of CRR
  Art.~77--78 is approximated by the MDA threshold.
  \item \textbf{No systemic or sovereign factor.} The spread process is issuer-specific and
  capital jumps do not move spreads. AT1 spreads widened by hundreds of basis points in March
  2020 and March 2023 with no change in capital, and BNP's spreads widened with French
  sovereign stress after the valuation date. A common sector and sovereign factor with
  co-jumps in capital and spreads is the main extension.
  \item \textbf{Rates.} The current curve is used for all future resets; a short-rate model
  would add rate-driven extension.
  \item \textbf{Unidentified spread parameters.} The spread volatility $\eta$ and the capital
  premium $\beta$ are assumed. A cross-section of issuers, or a regression of AT1 spreads on the
  distance to the MDA threshold, would calibrate $\beta$.
  \item \textbf{Data.} Prices are evaluated rather than traded; the sample covers fifteen
  months of relatively tight spreads; the bid--offer is of the same size as the cross-bond
  dispersion; the CDS is in EUR and the bonds in USD. TRACE transaction prices would allow a
  test on traded data.
  \item \textbf{Conversion.} A single stressed share price replaces a joint model of equity and
  capital; with conversion worth $4$bp this is immaterial for BNP, but not for an issuer close
  to its trigger.
\end{enumerate}

\subsection{Extensions}

Three extensions follow naturally. A common factor in spreads and capital, estimated on a
cross-section of European AT1 issuers, would capture the systemic component of the risk and
allow the AT1-specific premium to be decomposed into issuer and sector parts. A model of
distributions in which the bank allocates the MDA between dividends and coupons would
replace the upper bound on the MDA cost by an estimate. And a joint calibration across
issuers, with $\lres$ anchored to each issuer's Tier~2, would turn the residual premium into a
relative-value signal.

%======================================================================
\section{Conclusion}\label{sec:conclusion}
%======================================================================

AT1 bonds are claims on bank capital, and modelling them as such makes their three main
contingencies, the trigger, the coupon restriction and the call, consistent by construction.
Capital alone cannot explain their prices: the model needs a refinancing spread that drives
the call and a write-down hazard that rises in the same states. Conditioning on the state
paths makes the price an explicit, smooth and convex function of the hazard, so the implied
hazard is unique and every comparative static is exact on a given set of paths.

The implied hazard is not a black box. Because Tier~2 faces the same resolution event as AT1,
pricing the issuer's Tier~2 on the same state splits the hazard into a part Tier~2 also
prices and an AT1-specific premium. For BNP Paribas in September 2026, about two-thirds of
the spread, some $180$bp, is priced in Tier~2 as well. This part survives a reputational
call rule, every attribution that prices the hazard before the contract features, a term
structure of the resolution hazard and the range of capital dynamics that history and the
Tier~2 jointly admit. The remaining $89$--$98$bp pays for being an AT1 rather than a Tier~2, and is
equally robust. How it divides among expected MDA coupon cuts, extension and an
AT1-specific premium depends on how fast the bank rebuilds capital after a shock: the
premium is about $50$bp in the base case and zero at the historical estimate. The same dependence drives the level of the capital
stress. The Tier~2 narrows both, because the creditor hierarchy forbids pricing the AT1 as
safer than the Tier~2, but it does not remove them, and both are reported with their
range.